% -*- coding: utf-8 -*-
% 主论文主文件，基于 jnuthesis 模板，并将绪论拆分为独立文件。
\documentclass[14pt]{jnuthesis}
\usepackage{xeCJK}
\usepackage{amsmath,amssymb}
\usepackage{algorithm}
\usepackage{algpseudocode}
\usepackage{tikz}
\usepackage{pgfplots}
\usepackage{comment}
\pgfplotsset{compat=1.18}
\setlength{\headheight}{13pt}
% 使用类文件 jnuthesis.cls 中对章节样式的默认设置（宋体，三号，加粗，左对齐）
% 如需更改，请编辑 jnuthesis.cls 中的 CTEXsetup 对应项。

% 不强制章节在新页（取消双页清页），使章节尽量紧接排列
\let\cleardoublepage\relax

\renewcommand{\biaoti}{Null}
\renewcommand{\xueyuan}{网络空间安全学院}
\renewcommand{\zhuanye}{Null}
\renewcommand{\xingming}{Null}
\renewcommand{\xuehao}{000000000}
\renewcommand{\kecheng}{Null}
\renewcommand{\daoshi}{Null}

\begin{document}
\titlepage
\pagenumbering{arabic}
\tableofcontents

\addcontentsline{toc}{chapter}{\hspace{-1em}摘要}

\zhaiyao
empty

\guanjianci
empty

\newpage



\begin{comment}
{
\let\clearpage\relax
\let\cleardoublepage\relax
\input{chapters/chapter1_introduction}
\input{chapters/chapter2_prelim}
\input{chapters/chapter3_system_model}
\input{chapters/chapter4_cube}
\input{chapters/chapter5_security}
\input{chapters/chapter6_performance}
\input{chapters/chapter7_conclusion}
}
\end{comment}

\section{Test}
Test~\cite{boneh2004short}

\bibliographystyle{unsrt}\addcontentsline{toc}{chapter}{\hspace{-1em}参考文献}

\bibliography{reference}

\appendix

\end{document}
